\section{Limitations and Ethical Considerations}

This manuscript reports a completed development decision, not a formal-test result. The 400-claim feasibility packet starts from released judgment candidates and demonstrably overestimates performance relative to the label-blind full-corpus comparison. The development sample contains 367 eligible claims after the fixed exclusion policy rather than all released train/development claims. Expected-F1 decoding assumes conditionally independent evidence-membership variables even though documents and duplicated images may be dependent. Temperature scaling cannot correct a poorly ranked candidate pool, and the simple logistic scorer is not a capacity benchmark. Robustness tests, the five-seed formal matrix, and external validation were not run because the primary development gate failed. MOCHEG contains inherited professional fact-check labels and English-language web evidence; results may not transfer to other languages, domains, or live information environments. Public mirrors reproduce official counts but remain secondary transports, and uncertainty is claim-level because reliable event grouping is unavailable. These constraints prevent any claim of deployment readiness or general fact-checking reliability.

The system is intended to assist evidence review, not autonomously determine truth. Incorrect verdicts or incomplete evidence can misrepresent people, events, and sources. The project uses public research data, records provenance and immutable identifiers, and does not redistribute credentials or private user data. Generated explanations are excluded from the primary claim because fluent text is not evidence of correct attribution. Formal outputs must preserve probabilities, selected identifiers, hashes, and failure records so that unfavorable results remain auditable.
